\documentclass[12pt,a4paper]{report}
\usepackage[utf8]{inputenc}
\usepackage[T1]{fontenc}
\usepackage[french]{babel}
\usepackage{geometry}
\geometry{margin=2.5cm}
\usepackage{graphicx}
\usepackage{xcolor}
\usepackage{listings}
\usepackage{hyperref}
\usepackage{fancyhdr}
\usepackage{titlesec}
\usepackage{tcolorbox}
\usepackage{enumitem}
\usepackage{booktabs}
\usepackage{longtable}
\usepackage{tikz}
\usetikzlibrary{shapes.geometric, arrows.meta, positioning, fit}

\definecolor{primary}{HTML}{1a1a2e}
\definecolor{accent}{HTML}{16a34a}
\definecolor{codebg}{HTML}{f5f5f5}
\definecolor{codeframe}{HTML}{e0e0e0}

\hypersetup{colorlinks=true,linkcolor=primary,urlcolor=accent,citecolor=primary}

\lstset{
  backgroundcolor=\color{codebg},
  frame=single,
  rulecolor=\color{codeframe},
  basicstyle=\ttfamily\footnotesize,
  keywordstyle=\color{blue}\bfseries,
  stringstyle=\color{accent},
  commentstyle=\color{gray}\itshape,
  breaklines=true,
  showstringspaces=false,
  tabsize=2,
  numbers=left,
  numberstyle=\tiny\color{gray},
  xleftmargin=1.5em,
  framexleftmargin=1em,
}

\lstdefinelanguage{TypeScript}{
  keywords={import, export, from, const, let, var, function, return, if, else, class, new, async, await, try, catch, throw, interface, type, enum, extends, implements, private, public, protected, static, readonly, as, typeof, void, null, undefined, true, false, Promise},
  morecomment=[l]{//},
  morecomment=[s]{/*}{*/},
  morestring=[b]',
  morestring=[b]",
  morestring=[b]`,
  sensitive=true
}

\pagestyle{fancy}
\fancyhf{}
\fancyhead[L]{\footnotesize\textcolor{primary}{Propre \& Confiance -- Documentation Interne}}
\fancyhead[R]{\footnotesize\textcolor{primary}{\leftmark}}
\fancyfoot[C]{\thepage}
\renewcommand{\headrulewidth}{0.4pt}

\titleformat{\chapter}[display]
  {\normalfont\huge\bfseries\color{primary}}{\chaptertitlename\ \thechapter}{20pt}{\Huge}
\titleformat{\section}{\normalfont\Large\bfseries\color{primary}}{\thesection}{1em}{}
\titleformat{\subsection}{\normalfont\large\bfseries\color{primary!80}}{\thesubsection}{1em}{}

\newtcolorbox{infobox}[1][]{colback=accent!5,colframe=accent,fonttitle=\bfseries,title=#1}
\newtcolorbox{warnbox}[1][]{colback=red!5,colframe=red!70,fonttitle=\bfseries,title=#1}

\title{
  \vspace{-2cm}
  \textcolor{primary}{\rule{\linewidth}{2pt}}\\[0.5cm]
  {\Huge\textbf{\textcolor{primary}{Propre \& Confiance}}}\\[0.3cm]
  {\LARGE\textcolor{primary!70}{Documentation Technique d'Architecture Backend}}\\[0.5cm]
  \textcolor{primary}{\rule{\linewidth}{2pt}}\\[1.5cm]
  {\large\textcolor{gray}{Document interne -- Ne pas diffuser}}\\[0.3cm]
  {\normalsize Version 1.0 -- Mai 2026}
}
\author{
  \textbf{Équipe Backend -- Benflux Company}\\
  \texttt{benjamin@benflux.com}
}
\date{}

\begin{document}
\maketitle
\thispagestyle{empty}

\begin{warnbox}[DOCUMENT CONFIDENTIEL]
Ce document contient des détails sensibles sur l'architecture interne du système. Il ne doit en aucun cas être poussé sur GitHub ni partagé en dehors de l'équipe d'ingénierie.
\end{warnbox}

\clearpage
\tableofcontents
\clearpage

%% ======================================================================
\chapter{Introduction et Vision du Projet}
%% ======================================================================

\section{Contexte}
\textbf{Propre \& Confiance} est une plateforme de mise en relation entre des \textbf{familles} à la recherche de services de ménage et des \textbf{ménagères} professionnelles. Le backend expose une API RESTful qui gère l'ensemble du cycle de vie : inscription, vérification d'identité (KYC), réservation de missions, paiement, notation et administration.

\section{Stack Technologique}
\begin{longtable}{@{}p{4cm}p{9cm}@{}}
\toprule
\textbf{Composant} & \textbf{Technologie} \\
\midrule
Langage & TypeScript 5.5 (strict mode) \\
Runtime & Node.js avec ts-node-dev (dev) \\
Framework HTTP & Express.js 4.19 \\
ORM & TypeORM 0.3 \\
Base de données & PostgreSQL 16 \\
Authentification & JWT (Access + Refresh Tokens) \\
Validation & class-validator + class-transformer \\
Upload fichiers & Multer (mémoire) + Cloudinary \\
Email & Nodemailer (SMTP / Ethereal en dev) \\
Logging & Winston + DailyRotateFile \\
Documentation API & Swagger (swagger-jsdoc + swagger-ui-express) \\
Process Manager & PM2 (cluster mode en production) \\
Sécurité & Helmet, CORS, express-rate-limit, bcryptjs \\
\bottomrule
\end{longtable}

\section{Principes Architecturaux}
\begin{enumerate}
  \item \textbf{Modularité} -- Chaque domaine métier est isolé dans son propre module (Controller / Service / DTO / Routes).
  \item \textbf{Séparation des responsabilités} -- Les middlewares gèrent les préoccupations transversales (auth, validation, erreurs).
  \item \textbf{Sécurité en profondeur} -- Chaque couche applique ses propres contrôles (RBAC, rate-limiting, validation stricte).
  \item \textbf{Convention over Configuration} -- Nommage et structure uniformes dans tout le projet.
\end{enumerate}

%% ======================================================================
\chapter{Structure du Projet}
%% ======================================================================

\section{Arborescence des Dossiers}
\begin{lstlisting}[language={},title={Arborescence src/}]
src/
|-- server.ts              # Point d'entree (bootstrap)
|-- app.ts                 # Configuration Express
|-- config/                # Configurations centralisees
|   |-- database.config.ts
|   |-- jwt.config.ts
|   |-- logger.config.ts
|   +-- swagger.config.ts
|-- common/                # Code partage (transversal)
|   |-- enums/             # Enumerations metier
|   |-- types/             # Types TypeScript globaux
|   +-- utils/             # Fonctions utilitaires
|-- entities/              # Entites TypeORM (modeles DB)
|-- middleware/             # Middlewares Express
|-- modules/               # Modules metier
|   |-- auth/              # Authentification
|   |-- user/              # Gestion utilisateurs
|   |-- maid/              # Profils menageres
|   |-- kyc/               # Verification d'identite
|   |-- ticket/            # Missions / reservations
|   |-- payment/           # Paiements
|   |-- wallet/            # Portefeuilles
|   |-- withdrawal/        # Retraits
|   |-- rating/            # Notations
|   |-- search/            # Recherche
|   |-- notification/      # Notifications
|   +-- admin/             # Administration
|-- services/              # Services externes (mail)
|-- seeds/                 # Donnees initiales
+-- migrations/            # Migrations TypeORM
\end{lstlisting}

\section{Patron d'un Module}
Chaque module suit une structure identique :
\begin{lstlisting}[language={}]
modules/<nom>/
|-- <nom>.controller.ts    # Gere les requetes HTTP
|-- <nom>.service.ts       # Logique metier
|-- <nom>.routes.ts        # Definition des routes
+-- dto/                   # Data Transfer Objects
    +-- <action>-<nom>.dto.ts
\end{lstlisting}

\begin{infobox}[Pourquoi cette structure ?]
Ce découpage permet de : (1) tester chaque couche indépendamment, (2) remplacer une implémentation sans toucher aux autres couches, (3) naviguer intuitivement dans le code.
\end{infobox}

%% ======================================================================
\chapter{Cycle de Vie d'une Requête}
%% ======================================================================

\section{Pipeline de Middlewares}

\begin{center}
\begin{tikzpicture}[
  box/.style={rectangle, draw=primary, fill=primary!10, rounded corners, minimum width=3cm, minimum height=0.8cm, font=\small\bfseries, text=primary},
  arr/.style={-{Stealth[length=3mm]}, thick, primary},
  node distance=0.6cm
]
\node[box] (helmet) {Helmet};
\node[box, below=of helmet] (cors) {CORS};
\node[box, below=of cors] (body) {Body Parser};
\node[box, below=of body] (logger) {HTTP Logger};
\node[box, below=of logger] (rate) {Rate Limiter};
\node[box, below=of rate] (auth) {Authenticate};
\node[box, below=of auth] (email) {Email Verified};
\node[box, below=of email] (roles) {Authorize (RBAC)};
\node[box, below=of roles] (validate) {Validate Body};
\node[box, below=of validate, fill=accent!20] (ctrl) {Controller};

\draw[arr] (helmet) -- (cors);
\draw[arr] (cors) -- (body);
\draw[arr] (body) -- (logger);
\draw[arr] (logger) -- (rate);
\draw[arr] (rate) -- (auth);
\draw[arr] (auth) -- (email);
\draw[arr] (email) -- (roles);
\draw[arr] (roles) -- (validate);
\draw[arr] (validate) -- (ctrl);
\end{tikzpicture}
\end{center}

\section{Détail de chaque middleware}

\subsection{Helmet}
Ajoute des en-têtes HTTP de sécurité (CSP, HSTS, X-Frame-Options, etc.).

\subsection{CORS}
Contrôle les origines autorisées via \texttt{ALLOWED\_ORIGINS}. Supporte le mode wildcard pour le développement et les origines explicites pour la production. Le flag \texttt{credentials: true} permet l'envoi de cookies (refresh token).

\subsection{Rate Limiting}
Deux niveaux : un \textbf{limiteur global} (100 requêtes / 15 min) et un \textbf{limiteur d'authentification} (10 requêtes / 15 min) pour les routes sensibles (login, register, forgot-password).

\subsection{Authentification JWT}
\begin{lstlisting}[language=TypeScript,title={auth.middleware.ts -- Extrait}]
const payload = jwt.verify(token, jwtConfig.secret);
// Verification: mot de passe change apres emission du token
const user = await userRepo.findOne({
  where: { id: payload.sub },
  select: ['passwordChangedAt'],
});
if (user?.passwordChangedAt) {
  const changedTs = Math.floor(
    user.passwordChangedAt.getTime() / 1000
  );
  if (payload.iat < changedTs) {
    // Token invalide: mot de passe change
  }
}
req.user = { id: payload.sub, email, role };
\end{lstlisting}

\subsection{Autorisation RBAC}
\begin{lstlisting}[language=TypeScript,title={roles.middleware.ts}]
export function authorize(...allowedRoles: UserRole[]) {
  return (req, res, next) => {
    if (!allowedRoles.includes(req.user.role)) {
      return ResponseUtil.forbidden(res,
        'Access restricted. Required: '
        + allowedRoles.join(', '));
    }
    next();
  };
}
\end{lstlisting}

Trois rôles : \texttt{FAMILY}, \texttt{MAID}, \texttt{ADMIN}. Chaque route déclare explicitement les rôles autorisés.

\subsection{Validation des Données}
Utilise \texttt{class-validator} avec les options \texttt{whitelist: true} et \texttt{forbidNonWhitelisted: true} pour rejeter tout champ inconnu. Les DTOs définissent les règles de validation via des décorateurs.

%% ======================================================================
\chapter{Modèle de Données}
%% ======================================================================

\section{Diagramme Entité-Relation}

\begin{center}
\begin{tikzpicture}[
  entity/.style={rectangle, draw=primary, fill=primary!8, rounded corners=3pt, minimum width=2.8cm, minimum height=0.7cm, font=\small\bfseries},
  rel/.style={-{Stealth}, thick, accent!70},
  node distance=1.5cm and 2.5cm
]
\node[entity] (user) {User};
\node[entity, below left=of user] (family) {Family};
\node[entity, below=of user] (maid) {Maid};
\node[entity, below right=of user] (admin) {Admin};
\node[entity, right=2.5cm of maid] (kyc) {Kyc};
\node[entity, below=of family] (ticket) {Ticket};
\node[entity, below=of ticket] (payment) {Payment};
\node[entity, right=of ticket] (rating) {Rating};
\node[entity, below=of maid] (wallet) {Wallet};
\node[entity, right=of wallet] (withdrawal) {Withdrawal};
\node[entity, above right=1cm and 2.5cm of user] (notif) {Notification};

\draw[rel] (user) -- node[left,font=\tiny]{1:1} (family);
\draw[rel] (user) -- node[left,font=\tiny]{1:1} (maid);
\draw[rel] (user) -- node[right,font=\tiny]{1:1} (admin);
\draw[rel] (maid) -- node[above,font=\tiny]{1:1} (kyc);
\draw[rel] (family) -- node[left,font=\tiny]{1:N} (ticket);
\draw[rel] (maid) -- node[right,font=\tiny]{1:N} (ticket);
\draw[rel] (ticket) -- node[left,font=\tiny]{1:1} (payment);
\draw[rel] (ticket) -- node[above,font=\tiny]{1:N} (rating);
\draw[rel] (maid) -- node[left,font=\tiny]{1:1} (wallet);
\draw[rel] (maid) -- node[right,font=\tiny]{1:N} (withdrawal);
\draw[rel] (user) -- node[above,font=\tiny]{1:N} (notif);
\end{tikzpicture}
\end{center}

\section{Entités Principales}

\subsection{User}
Entité centrale. Contient les champs d'authentification (\texttt{email}, \texttt{password}, \texttt{role}), les flags de sécurité (\texttt{isEmailVerified}, \texttt{isActive}), et les tokens sensibles marqués \texttt{select: false} pour ne jamais les exposer par défaut.

\subsection{Family / Maid / Admin}
Profils spécialisés liés en \texttt{OneToOne} à \texttt{User}. Cette approche de \textbf{Table-per-Type} évite les colonnes nulles dans une table unique tout en permettant des champs spécifiques (ex: \texttt{hourlyRate} pour Maid, \texttt{commissionRate} pour Admin).

\subsection{Ticket}
Représente une mission de ménage. Contient le calcul de prix complet : \texttt{subtotal}, \texttt{discountRate}, \texttt{discountAmount}, \texttt{totalPrice}, \texttt{commissionRate}. Cycle de vie : \texttt{PENDING} $\to$ \texttt{PAID} $\to$ \texttt{IN\_PROGRESS} $\to$ \texttt{COMPLETED}.

\subsection{Payment}
Lié 1:1 au Ticket. Supporte la validation manuelle par l'admin (preuve de paiement via screenshot). Champs d'audit : \texttt{reviewedBy}, \texttt{reviewedAt}, \texttt{distributedAt}, \texttt{escrowedAt}.

\subsection{Wallet}
Portefeuille polymorphe : peut appartenir à une Family, une Maid ou un Admin via des clés étrangères nullables (\texttt{maidId}, \texttt{adminId}, \texttt{familyId}).

%% ======================================================================
\chapter{Sécurité}
%% ======================================================================

\section{Authentification JWT Dual-Token}

\begin{infobox}[Stratégie Access + Refresh]
\begin{itemize}[nosep]
  \item \textbf{Access Token} : durée courte (15 min), envoyé dans le header \texttt{Authorization: Bearer}.
  \item \textbf{Refresh Token} : durée longue (7 jours), stocké en base de données (hashé) et en cookie HttpOnly.
  \item Rotation automatique : chaque \texttt{/auth/refresh} génère une nouvelle paire et invalide l'ancien refresh token.
\end{itemize}
\end{infobox}

\section{Hachage des Mots de Passe}
Utilisation de \texttt{bcryptjs} avec un coût de \textbf{12 rounds}. Le champ \texttt{password} est marqué \texttt{select: false} dans l'entité TypeORM pour ne jamais être retourné par défaut.

\section{Vérification d'Email par OTP}
\begin{enumerate}
  \item À l'inscription, un OTP à 6 chiffres est généré via \texttt{crypto.randomInt}.
  \item L'OTP est hashé en SHA-256 avant stockage en base (\texttt{emailVerificationToken}).
  \item L'utilisateur soumet l'OTP, le serveur hash la soumission et compare.
  \item Expiration : 24 heures.
\end{enumerate}

\section{Protection contre les Attaques}
\begin{longtable}{@{}p{4cm}p{9cm}@{}}
\toprule
\textbf{Attaque} & \textbf{Protection} \\
\midrule
Brute force & Rate limiting (10 req/15 min sur auth) \\
XSS & Helmet CSP headers \\
CSRF & Cookies SameSite, CORS strict \\
Injection SQL & ORM paramétré (TypeORM) \\
Token replay & Vérification \texttt{passwordChangedAt} vs \texttt{iat} \\
Mass assignment & \texttt{whitelist: true, forbidNonWhitelisted: true} \\
Data leakage & \texttt{select: false} sur champs sensibles + \texttt{sanitizeUser()} \\
\bottomrule
\end{longtable}

%% ======================================================================
\chapter{Logique Métier}
%% ======================================================================

\section{Flux de Réservation (Ticket)}
\begin{enumerate}
  \item La famille crée un ticket en spécifiant : ménagère, tâches, durée, heures/jour.
  \item Le système calcule le prix via \texttt{calculateMissionPrice()} avec remise progressive.
  \item La famille soumet une preuve de paiement (screenshot Airtel Money).
  \item L'admin valide manuellement $\to$ statut passe à \texttt{PAID}.
  \item Après complétion : commission prélevée, fonds distribués vers les wallets.
\end{enumerate}

\section{Grille de Remises}
\begin{lstlisting}[language=TypeScript,title={pricing.util.ts}]
// 1-2 jours  -> 0%    | 3-6 jours  -> 5%
// 7-13 jours -> 10%   | 14-29 jours -> 15%
// 30+ jours  -> 20%
function getDiscountRate(days: number): number {
  if (days >= 30) return 20;
  if (days >= 14) return 15;
  if (days >= 7) return 10;
  if (days >= 3) return 5;
  return 0;
}
\end{lstlisting}

\section{Split de Paiement}
\begin{lstlisting}[language=TypeScript]
function splitPayment(totalPrice, commissionRate) {
  const adminAmount = totalPrice * (commissionRate / 100);
  const maidAmount = totalPrice - adminAmount;
  return { adminAmount, maidAmount };
}
\end{lstlisting}

\section{Module KYC}
Processus de vérification d'identité pour les ménagères. Documents requis : pièce d'identité (recto/verso), selfie, casier judiciaire. Upload parallèle vers Cloudinary via \texttt{Promise.all}. Cycle : \texttt{PENDING} $\to$ \texttt{APPROVED} / \texttt{REJECTED}.

%% ======================================================================
\chapter{Services Externes}
%% ======================================================================

\section{Cloudinary (Stockage Fichiers)}
Les fichiers sont stockés en mémoire (Multer \texttt{memoryStorage}) puis envoyés vers Cloudinary via un stream. Organisation des dossiers : \texttt{propre-confiance/kyc/<userId>/}.

\section{Nodemailer (Emails)}
Templates HTML intégrés pour les OTP (vérification email et reset mot de passe). En développement : Ethereal Mail (bac à sable). Configuration flexible avec fallback \texttt{SMTP\_*} / \texttt{MAIL\_*}.

\section{Winston (Logging)}
\begin{itemize}[nosep]
  \item \textbf{Développement} : console colorisée, niveau \texttt{debug}.
  \item \textbf{Production} : fichiers rotatifs journaliers (\texttt{logs/app-YYYY-MM-DD.log}), compressés, rétention 14 jours, taille max 20 Mo.
\end{itemize}

%% ======================================================================
\chapter{Scalabilité et Production}
%% ======================================================================

\section{PM2 Cluster Mode}
\begin{lstlisting}[language={},title={ecosystem.config.cjs}]
{
  name: 'propre-confiance-api',
  script: 'dist/server.js',
  instances: 'max',       // 1 worker par CPU
  exec_mode: 'cluster',
  max_memory_restart: '500M',
}
\end{lstlisting}

\section{Graceful Shutdown}
Le serveur intercepte \texttt{SIGTERM} et \texttt{SIGINT}, ferme proprement les connexions HTTP, puis quitte. Un timeout de 10 secondes force l'arrêt si le serveur ne répond pas.

\section{Pagination}
Toutes les listes sont paginées avec métadonnées : \texttt{total}, \texttt{page}, \texttt{limit}, \texttt{totalPages}, \texttt{hasNextPage}, \texttt{hasPrevPage}. Limite maximale : 100 éléments par page.

\section{Réponse Standardisée}
\begin{lstlisting}[language=TypeScript]
// Succes
{ success: true, message: "...", data: {...}, meta: {...} }
// Erreur
{ success: false, message: "...", errors: [...] }
\end{lstlisting}

%% ======================================================================
\chapter{Guide de Reproduction}
%% ======================================================================

Pour reproduire cette architecture dans un nouveau projet :

\begin{enumerate}
  \item \textbf{Initialiser} : \texttt{npm init}, installer Express, TypeORM, TypeScript.
  \item \textbf{Configurer TypeScript} : activer \texttt{strict}, \texttt{experimentalDecorators}, \texttt{emitDecoratorMetadata}.
  \item \textbf{Créer la structure} : \texttt{src/\{config, common, entities, middleware, modules, services\}}.
  \item \textbf{Configurer la DB} : créer \texttt{database.config.ts} avec \texttt{DataSource}.
  \item \textbf{Créer les entités} : une classe par table, décorateurs TypeORM.
  \item \textbf{Créer les middlewares} : error handler global, auth, RBAC, validation.
  \item \textbf{Créer les modules} : pour chaque domaine : routes $\to$ controller $\to$ service.
  \item \textbf{Assembler dans app.ts} : middlewares globaux puis routes.
  \item \textbf{Bootstrap dans server.ts} : init DB $\to$ create app $\to$ listen.
\end{enumerate}

\begin{infobox}[Règle d'or]
Un module ne doit \textbf{jamais} importer directement un autre module. La communication se fait via les services, et les dépendances transversales passent par \texttt{common/}.
\end{infobox}

%% ======================================================================
\chapter{Bonnes Pratiques}
%% ======================================================================

\begin{enumerate}
  \item \textbf{Ne jamais exposer les champs sensibles} : utiliser \texttt{select: false} et \texttt{sanitizeUser()}.
  \item \textbf{Valider toutes les entrées} : DTOs avec \texttt{class-validator}, \texttt{whitelist: true}.
  \item \textbf{Utiliser des enums} pour les statuts et rôles (pas de strings magiques).
  \item \textbf{Gérer les erreurs avec une hiérarchie} : \texttt{AppError} $\to$ \texttt{NotFoundError}, \texttt{UnauthorizedError}, etc.
  \item \textbf{Logger les événements importants} : authentifications, paiements, erreurs.
  \item \textbf{Versionner l'API} : préfixe \texttt{/api/v1/} pour permettre des évolutions futures.
  \item \textbf{Variables d'environnement} : jamais de secrets en dur, tout dans \texttt{.env}.
  \item \textbf{Graceful shutdown} : fermer proprement les connexions DB et HTTP.
  \item \textbf{Git workflow} : une branche par fonctionnalité, PR obligatoire avant merge.
  \item \textbf{Documentation Swagger} : synchronisée avec le code source.
\end{enumerate}

\vfill
\begin{center}
\textcolor{primary}{\rule{0.5\linewidth}{1pt}}\\[0.5cm]
{\large\textcolor{primary}{\textbf{Propre \& Confiance}}}\\
{\small\textcolor{gray}{Documentation générée le \today{} -- Benflux Company}}
\end{center}

\end{document}
